\myaddRA{\subsection{Instantiation III: Conflict-State DPCM for General \texorpdfstring{$r$}{r}}}

% For $r \ge 3$, updates operate on higher-order structures. 

In this subsection, we generalize the update strategy from vertices and edges to arbitrary $r$-cliques.
%
Under Theorem~\ref{thm:dpcm-states}, the sufficient state of an affected path now consists of its current \hold/\pivot partition together with the family of removed $r$-cliques that still conflict with the path.
%
Thus the general regime becomes a \emph{conflict-state DPCM}: each affected path maintains enough information to determine whether it survives, whether it must split, and how the surviving subpaths contribute to the remaining $r$-cliques.


% For $r\ge 3$, peeling removes the $r$-clique with minimum $s$-support, i.e., it eliminates all $s$-cliques that contain that $r$-clique. To realize this operation on the Succinct Clique Tree(\cpi) without enumerating $r$-cliques, we translate the deletion of an $r$-clique that lies on a single clique path into a split of that path into subpaths that avoid the clique.


% defined conflict set which is all verteicse in the clique to be removed -> 不能出现在同一个clique path
% while in the meantime retain all other cliques uniquely represented by one path
% each vertex in the clique can be either hold or pivot, when edge (r=2), three cases, r arbitrary number, r+1 possible cases
% we may split the path into multiple paths 
% we observed that the number of new path(s) depends on the number of pivot vertice in the clique




% \stitle{From Edges to Cliques.}
% %
% For $r \ge 3$, updates operate on higher-order structures, where removing an $r$-clique $R$ simultaneously affects all cliques containing it. 

% \stitle{Conflict Sets and Path Validity.}
% Updates operate on higher-order structures, where removing an $r$-clique $R$ simultaneously affects all cliques containing it. 

% We define the \emph{conflict set} of a removed $r$-clique $R$ as the set of its vertices; after the update, no path in the \cpi may contain all vertices in this set, while all other cliques remain uniquely represented. 
% %
% Each vertex in $R$ serves as a \hold or \pivot within a path, determining the update behavior. 
% % For $r=2$, we observe three cases (\hold{–}\hold, \hold{–}\pivot, \pivot{–}\pivot); in general, $(r{+}1)$ combinations exist. 
% %
% Removing an $r$-clique thus requires resolving these combinations and, when necessary, splitting affected paths to maintain correctness. 
% The number of new paths generated depends on the number of \pivot vertices in $R$. 
% This insight underpins our generalized update algorithm, which efficiently supports arbitrary $r$. 
% %
% Building on this idea, we extend the \textit{path-splitting} operation to support $r$-clique removals for arbitrary $r$, as formalized in Theorem~\ref{thm:split}.

\stitle{Conflict Sets and Path Validity.}
Removing an $r$-clique $R$ is a higher-order update.
It simultaneously affects all $s$-cliques that contain $R$.
% 
We define the \emph{conflict set} of $R$ as its vertex set $V(R)$.
After the update, no clique path may realize all vertices in $V(R)$ at the same time.
Meanwhile, every remaining $s$-clique must still be encoded by exactly one path.
% 
On a path, vertices of $R$ can appear as \hold or \pivot.
Holds are mandatory, while pivots represent alternative choices.
Therefore, updating the CPI reduces to removing exactly those realizations that include the whole conflict set.
If needed, we split the affected path into a small number of new paths to restore validity and uniqueness (Lemma~\ref{lem:clique_count}).
Crucially, the number of new paths depends only on how many vertices of $R$ act as pivots on that path.
Building on this observation, Theorem~\ref{thm:split} formalizes the general path-splitting rule for arbitrary $r$.

% Our goal is to update the CPI after removing an $r$-clique $R$,
% % 
% Updates operate on higher-order structures is challenging, removing an $r$-clique $R$ simultaneously affects \emph{all} cliques that contain it, and thus cannot be handled by local edge-style updates.

% We define the \emph{conflict set} of the removed $r$-clique $R$ as its vertex set.
% After the update, no clique path in the \cpi may contain \emph{all} vertices in this conflict set; meanwhile, every remaining clique must still be represented \emph{uniquely} by exactly one path.
% %
% Within a clique path, vertices of $R$ may appear as \hold or \pivot, which determines how the path reacts to the removal:
% holds are mandatory on the path, while pivots represent alternative choices.
% Therefore, removing $R$ amounts to eliminating exactly those path realizations that include the entire conflict set, while keeping the others intact.
% %
% \myadd{
% To restore validity and uniqueness(Lemma \ref{lem:clique_co bnmunt}), we may need to split an affected path into a small number of new paths that exclude $R$.
% Crucially, the number of new paths depends only on how many vertices of $R$ act as pivots on that path.
% %
% Building on this idea, we generalize the \textit{path-splitting} operation to handle $r$-clique removals for arbitrary $r$, as formalized in Theorem~\ref{thm:split}.
% }

% For $r=2$, deleting $\{u,v\}$ on one \cpi path has three outcomes: (\hold--\hold) removes the path; (\hold--\pivot) drops the pivot endpoint; (\pivot--\pivot) splits the path into two branches—one forbids $v$, the other mandates $v$ and forbids $u$\, (Theorem~\ref{thm:remove_edge}). 
% 我们可以换种方法理解, 对于(\pivot--\pivot)的情况, 我们可以把它看作是由两个冲突点引起的路径分割. 给定一个路径$\TPath$和一个包含连个点的conflict set $C=\{u,v\}$, 我们需要将路径$\TPath$分割成多个子路径, 使得每个子路径都不包含这个冲突集$C$, 并且还要确保所有不包含冲突集$C$的clique都被这些子路径唯一的表示出来.
% We can reinterpret the (\pivot--\pivot) case as a path split induced by two conflict points. Given a path $\TPath$ and a conflict set $C=\{u,v\}$ containing two vertices, we need to split the path $\TPath$ into multiple subpaths such that each subpath excludes the conflict set $C$, while ensuring that all cliques not containing $C$ are uniquely represented by these subpaths.


% formalizes this construction and proves its correctness.

% This is the $|C|=2$ instance of a more general split by a \emph{conflict set}.

% For a path $\TPath$ and an $r$-clique $C\subseteq \TPath$, order the pivot vertices of $C$ as $v_0,\ldots,v_{t-1}$. We produce $t$ subpaths that (i) promote $v_0,\ldots,v_{i-1}$ to \holds\ and (ii) exclude $v_i$. Theorem~\ref{thm:split} formalizes this construction and proves that the resulting subpaths are valid, encode no clique containing $C$, and together encode exactly the cliques on $\TPath$ that survive the deletion of $C$.
% Short/duplicated statement and sketch of Theorem~\ref{thm:split} removed here to avoid conflicting/ambiguous macro usage;
% the full and braced version of Theorem~\ref{thm:split} appears below.


% \begin{theorem}[Path splitting by one conflict set]\label{thm:split}
% Let $\TPath=(V_p^{\TPath},V_h^{\TPath})$ be a clique path in the \cpi\ of $G$, and let $C=(V_p^C,V_h^C)$ be an $r$-clique with $V_p^C\subseteq V_p^{\TPath}$ and $V_h^C\subseteq V_h^{\TPath}$. Write $V_p^C=\{v_0,\ldots,v_{t-1}\}$ with $t:=|V_p^C|\ge 1$. After removing all cliques that contain $C$, $\TPath$ is replaced by the following set of subpaths:
% \[
%   \mathcal{T}'=\{\,\TPath^{\mathrm{sub}}_i=(V_{p,i},V_{h,i}) : i=0,\ldots,t-1\,\},
% \]
% where
% \[
%   V_{h,i}= V_h^{\TPath} \cup \{v_0,\ldots,v_{i-1}\},\qquad
%   V_{p,i}= V_p^{\TPath} \setminus \{v_0,\ldots,v_i\},
% \]
% For all $s$-cliques encoded by $\TPath$ that do not contain $C$, exactly one subpath in $\mathcal{T}'$ encodes it.
% For all $s$-cliques $C_s$ encoded by $\TPath$ that $C \not\subseteq C_s$, exactly one subpath in $\mathcal{T}'$ encodes it.
% For a fixed size $k$, discard any subpath $\TPath^{\mathrm{sub}}_i$ that is infeasible for $k$ (Lemma~\ref{lem:path_pruning}); the number of size-$k$ cliques on any remaining subpath is given by Lemma~\ref{lem:clique_count}.
% \end{theorem}

\begin{theorem}[Clique Removal]\label{thm:split}

Let $\TPath = (V_p^{\TPath}, V_h^{\TPath})$ be a clique path, and let $C = (V_p^{C}, V_h^{C})$ be an $r$-clique such that $V_p^{C} \subseteq V_p^{\TPath}$ and $V_h^{C} \subseteq V_h^{\TPath}$. 
Given an ordered set $V_p^{C} = \{v_0, \ldots, v_{t-1}\}$, the updated path set $\TPath'$ is defined as:
\[
  \TPath' = \bigl\{\, (V_p^{\TPath} \setminus \{v_0, \ldots, v_i\},\; V_h^{\TPath} \cup \{v_0, \ldots, v_{i-1}\}) \;\big|\; i = 0, \ldots, t-1 \,\bigr\}.
\]
% Then, every $s$-clique $C_s$ encoded by $\TPath$ with $C \nsubseteq C_s$ is encoded by exactly one subpath in $\TPath'$, and no subpath encodes a clique containing $C$.
\end{theorem}


\begin{proof}
On a \cpi\ path, holds are mandatory and pivots optional. For each $i$, promoting $v_0,\ldots,v_{i-1}$ to holds and deleting $v_i$ keeps the path valid and forbids any clique that contains all pivots of $C$, so no subpath encodes $C$ (soundness). For any surviving clique $K$, let $i$ be the first index with $v_i\notin K$; then $K$ contains the promoted prefix and excludes $v_i$, hence is encoded precisely on subpath $i$ and on no other (completeness \& uniqueness).
\end{proof}

\begin{example}
\label{ex:split_example}

Assume we have a clique path $\TPath = \langle v_1, v_2, v_3, v_4, v_5, v_6 \rangle$,
where $V_p = \{v_1, v_2, v_3, v_6\}$ and $V_h = \{v_4, v_5\}$.
If we have an $r$-clique $C = \{v_1, v_2, v_3, v_4\}$, 
then $V_p^C = \{v_1, v_2, v_3\}$ and $V_h^C = \{v_4\}$.
According to Theorem~\ref{thm:split}, we can split $\TPath$ into three subpaths:
$V_{p,0} = \{v_1, v_2,v_6\},  V_{h,0} = \{v_4,v_5\},$
$V_{p,1} = \{v_1,v_6\}, V_{h,1} = \{v_3,v_4,v_5\},$
$V_{p,2} = \{v_6\}, V_{h,2} = \{v_2,v_3,v_4,v_5\}.$

\end{example}

\begin{remark}
Theorem~\ref{thm:split} allows a clique path $\TPath$ in the \cpi to be split into multiple paths that exclude the $r$-clique $C$. 
This operation is crucial, as it implicitly removes $r$-cliques—without explicit enumeration—while preserving the correctness of the remaining structure.
\end{remark}

\stitle{Role in DPCM.}
Theorem~\ref{thm:split} is the fundamental local state-transition operator for the general regime.
%
It is exact but potentially expensive when many conflicting cliques hit the same path.
%
The next subsection therefore adds a stronger journal-level refinement: dead-path-aware maintenance, which uses the maintained conflict state to certify that many affected paths can be discarded without recursive splitting.

% 基于此, 我们推出了算法~\ref{alg:r3}来进行对于$r\geq 3$的情况的核心分解.
% Based on this, we propose Algorithm~\ref{alg:peeling_r3} for the core decomposition in the case of $r \geq 3$.

% \paragraph{Algorithm Description.}
% 与Algorithm~\ref{alg:r1}和Algorithm~\ref{alg:r2}类似,我们先计算每个r-clqiue的s-support, 然后删除s-support最小的r-clique, 并更新受影响的r-clique的s-support, 直到所有的r-clique都被处理完.
% Similar to Algorithm~\ref{alg:peeling_r1} and Algorithm~\ref{alg:peeling_r2}, we first compute the s-support of each r-clique, then remove the r-clique with the smallest s-support, and update the s-support of the affected r-cliques, until all r-cliques have been processed.



\begin{algorithm}[t]
\small
\caption{\textsc{UpdateIndex}}
\label{alg:removeRClique}
\KwIn{clique path $\TPath{} = \{V_{H},V_{P}\}$, set of $r$-cliques to remove $C_{\mathrm{rm}}=\{C_0,\dots,C_{t-1}\}$, clique size $s$}
% \KwOut{All clique subpath produced after deleting $C_{\mathrm{rm}}$}
\KwOut{All subpaths induced after deleting $C_{\mathrm{rm}}$}
\SetKwFunction{Split}{PathSplit}
\SetKwProg{Fn}{Function}{}{}

% $V_{H} \gets$ hold vertices on $\TPath{}$;
% $V_{P} \gets$ pivot vertices on $\TPath{}$\\
Build inverted index $I(v) \gets \{\,i \mid v \in C_i\,\}$ for all $v \in \TPath$\\


cnt$\gets$ array of size $t$ initialized to $r$\;

\Split{$V_{H},\;V_{P}$}\;

\Fn{\Split{$V_{H},\;V_{P}$}}{

  \lIf{$|V_{H} \cup V_{P}| < s$ or $|V_{H}| > s$}{\Return} % pruning rule

  $C_j \gets$ the first $r$-clique in $C_{\mathrm{rm}}$ with $cnt[j] = r$\;
  
  \If{no such $C_j$ exists}{
     \textbf{output path }$\{V_{H}, V_{P}\}$\;
     
     \Return
  }

  \ForEach{$v \in C_j \cap V_{P}$}{
    \ForEach{$g \in I(v)$}{$cnt[g] \gets cnt[g]-1$}

    $V_{P} \gets V_{P} \setminus \{v\}$;

    \Split{$V_{H},\;V_{P}$}\;

    $V_{H} \gets V_{H} \cup \{v\}$\;

    \ForEach{$g \in I(v)$}{$cnt[g] \gets cnt[g]+1$}
  }
    
}

\end{algorithm}




% a vertex may appear in many cliques 
% we grounp the updates by vertices
% we build an inverted index for mapping vertex->clique
% update each vertex sequntially  (all cliques containing the vertex) instead of  each clique sequntially

\stitle{Batched Updates.}
When multiple cliques are removed simultaneously, they often overlap, leading to redundant updates on shared
vertices and paths. To mitigate this overhead, we aggregate overlap
ping updates into vertex-centric batches. Specifically, we construct
an $inverted$ $index$ that maps each vertex to the set of cliques con
taining it. Instead of processing cliques sequentially, we iterate over
vertices in this index and update all affected clique paths per ver
tex. This vertex-grouped strategy effectively merges overlapping
operations and reduce redundant updates.

% Theorem~\ref{thm:split} shows how deleting a single r-clique on one path induces a structured split. In peeling, however, a single path may need to remove multiple r-cliques at once. We therefore handle them in one pass: sweep the incident family once, branch only on pivots, and commit earlier branch choices as \holds. This yields exactly the subpaths that repeated single-clique splits would produce, while avoiding redundant splits and rebuilds (Algorithm~\ref{alg:removeRClique}).

\stitle{Algorithm.} 
Algorithm~\ref{alg:removeRClique} maintains a counter $cnt$ to track the number of vertices from each $r$-clique in $C_{\mathrm{rm}}$ that are still present on the current path. When $cnt$ equals $r$, it indicates that the path violates the conflict clique and requires splitting. The splitting process (Function \Split) recursively handles each $r$-clique to be removed.
First, in Line~6, we prune the path based on Lemma~\ref{lem:path_pruning}; if the current path cannot contain any cliques of size $s$, we return immediately.
Then, in Line~7, we select the first $r$-clique $C_j$ that needs to be removed.
If no such $r$-clique exists, it indicates that the current path does not require any deletion operations, so we report the current path as a subpath (Line~8).
Otherwise, we iterate through each vertex $v$ in the $r$-clique $C_j$, remove it from the path, and decrement the $cnt$ of all $r$-cliques containing $v$ by 1, as one vertex has been removed from these cliques.
In Lines 15-17, we promote all pivot nodes before $v$ to hold nodes according to Theorem~\ref{thm:split}, remove $v$ from the path, and then recursively call the function $\Split$ to process the updated path.
Finally, in lines 19-20, we remove $v$ from the path and continue processing the next vertex.
According to Theorem~\ref{thm:split}, we can ensure that after each path split, the resulting subpaths do not contain the removed $r$-clique, and all size-$s$ cliques that do not contain the removed $r$-clique are uniquely represented.

% \noindent\emph{Implementation note.}
% We represent vertex sets on clique paths (e.g., holds/pivots and conflict sets) using bitsets, enabling fast membership and set operations; this optimization applies starting from $r=3$ and is used throughout our implementation.


% \begin{theorem}[Correctness of Algorithm~\ref{alg:removeRClique}]
% \label{thm:removeRClique-correct}
% Let $\TPath=\{V_p^{\TPath},V_h^{\TPath}\}$ be a clique path of the \cpi and
% $\mathcal C=\{C_0,\dots,C_{t-1}\}$ a family of $r$-cliques contained in $\TPath$.
% Algorithm~\ref{alg:removeRClique} reports a family of clique subpath denoted as $\TPath^{Out} = [\TPath^{Out}_0,\dots,\TPath^{Out}_k]$, each represented as a pair $\{V_{Hold},V_{Piv}\}$.\
% The subpaths in $\TPath^{Out}$ satisfy the following properties:
% \begin{enumerate}
%   \item \textbf{Soundness.} Every path in $\TPath^{Out}$, does not cover any $r$-clique $C_i\in\mathcal C$.
%   \item \textbf{Completeness and uniqueness.} Every $s$-clique $C_s$ encoded by $\TPath$ that $C_i \not\subseteq C_s$ for all $C_i\in\mathcal C$ is encoded by exactly one path in $\TPath^{Out}$.
% \end{enumerate}
% Consequently, replacing $\TPath$ by the reported subpaths removes precisely the size-$s$ cliques that contain a member of $\mathcal C$.
% \end{theorem}
% \begin{proof}
% \emph{Soundness.}
% For each $C_i\in\mathcal C$, let $V_p^{C_i}$ be its pivot set.
% In each invocation of $\Split$, we either report the current path (line 8) or split it by the pivot set of some $C_j\in\mathcal C$ (line 7).
% In the former case, no $C_i$ can be contained in the reported path, since otherwise we would not have reported it.
% In the latter case, Theorem~\ref{thm:split} guarantees that none of the resulting subpaths contains $C_j$, and hence none contains any $C_i\in\mathcal C$.

% \emph{Completeness and uniqueness.}
% Let $C_s$ be any size-$s$ clique encoded by $\TPath$ that does not contain any $C_i\in\mathcal C$.
% We show that $C_s$ is encoded by exactly one path in $\TPath^{Out}$.
% Since $C_s$ does not contain any $C_i\in\mathcal C$, it also does not contain the pivot set $V_p^{C_i}$ of any $C_i$.
% Let $S_i:=C_s\cap V_p^{C_i}$ be the subset of conflict pivots contained in $C_s$.
% Since $S_i\neq V_p^{C_i}$, define $j_i:=\min\{j\in\{0,\dots,|V_p^{C_i}|-1\}\mid v_j\notin S_i\}$, i.e., the \emph{first missing} conflict pivot of $C_i$.
% Then $C_s$ contains $v_0,\dots,v_{j_i-1}$ and excludes $v_{j_i}$.
% In the execution of $\Split$, consider the first invocation that splits by $C_i$ (line 7).
% At that time, the current path still contains all vertices of $C_s$ (since no prior split could have removed any of them), so $C_s$ is encoded by that path.
% By Theorem~\ref{thm:split}, $C_s$ is then encoded by exactly one of the resulting subpaths, namely the one that makes $v_0,\dots,v_{j_i-1}$ holds and removes $v_{j_i}$.
% In subsequent invocations of $\Split$, $C_s$ remains encoded by the current path (since no split can remove any of its vertices), and by Theorem~\ref{thm:split}, it remains encoded by exactly one of the resulting subpaths.
% Continuing this argument for all $C_i\in\mathcal C$ shows that $C_s$ is encoded by exactly one path in $\TPath^{Out}$.
% \end{proof}

\stitle{Correctness.} 
We prove the correctness of Algorithm~\ref{alg:removeRClique} in terms of soundness, completeness, and uniqueness.
%
\emph{(Soundness)} 
The algorithm maintains a counter array $cnt$ to track how many vertices of each removed $r$-clique remain on the current path.
A path is emitted only when all $cnt$ values are strictly less than $r$, ensuring that the new path encodes no $r$-cliques marked for removal.
%
\emph{(Completeness \& Uniqueness)} 
When removing vertices of an $r$-clique iteratively, at each deletion step, all remaining cliques not containing the removed vertex are re-encoded uniquely in new paths (by Theorem~\ref{thm:split}). 
Since every vertex is removed exactly once, all valid combinations of remaining vertices are exhaustively enumerated and retained. 
Therefore, every surviving clique is preserved and encoded by exactly one path in the updated \cpi, ensuring both completeness and uniqueness.

\stitle{Implementation.}
The size of the residual graph is bounded by the maximum clique size of the input graph, which is typically small for real-world datasets (Table~\ref{tab:graph-stats}). 
This allows us to represent the residual graph as an adjacency matrix and employ \emph{bitsets} to encode vertex sets on paths, enabling highly efficient set operations (\textsc{and}, \textsc{or}, \textsc{xor}, \textsc{count}) through word-level parallelism~\cite{bitsetClique}.

% Every size-$s$ clique $C_s$ encoded by the original path $\TPath$ that does not contain any $r$-clique from $C_{\mathrm{rm}}$ is encoded by exactly some output paths.

% For losslessness, 
% % 算法每次递归只删除当前conflict clique $C_j$ 中的一个顶点$v$. 也就意味着在当前路径上, 增加任意一个顶点$v$都会导致有一个conflict clique被完全包含在路径上.
% The algorithm recursively removes one vertex $v$ from the current conflicting clique $C_j$ at a time. 
% This implies that adding any vertex $v$ to the current path would result in at least one conflicting clique being fully contained within the path.

% \emph{Uniqueness}: Each size-$s$ clique $C_s$ is encoded by exactly one output path.

% % does not contain any conflicting $r$-cliques, allowing us to output it as a subpath.
% % 
% % For losslessness, 假设有一个size-$s$ clique $C_s$ encoded by the original path $\TPath$ that does not包含任何来自$C_{\mathrm{rm}}$的$r$-clique.

% For losslessness, 
% % 算法每次递归只删除当前conflict clique $C_j$ 中的一个顶点$v$. 也就意味着在当前路径上, 增加任意一个顶点$v$都会导致有一个conflict clique被完全包含在路径上.
% The algorithm recursively removes one vertex $v$ from the current conflicting clique $C_j$ at a time. 
% This implies that adding any vertex $v$ to the current path would result in at least one conflicting clique being fully contained within the path.
% % 
% % 算法对于每个 conflicting clique会进行分割, 并且会尝试所有可能的分割方式.
% %  因此, 对于任意一个不包含来自$C_{\mathrm{rm}}$的$r$-clique的size-$s$ clique $C_s$, 它一定会在某次递归中被保留在当前路径上, 并且不会被删除的顶点$v$所影响.
% The algorithm splits the path for each conflicting clique and explores all possible splitting methods.
% Therefore, for any size-$s$ clique $C_s$ that does not contain any $r$-clique from $C_{\mathrm{rm}}$, it will inevitably be preserved on the current path during some recursion and will not be affected by the removed vertex $v$.
% % 
% % 并且根据Theorem~\ref{thm:split}, 它会被唯一的一个子路径所表示.

% For uniqueness, according to Theorem~\ref{thm:split}, each $s$-clique $C_s$ is represented by exactly one subpath. The algorithm strictly follows Theorem~\ref{thm:split} for each path split, ensuring that each size-$s$ clique $C_s$ is represented by exactly one subpath.

% 那么也就意味着

% \paragraph{Complexity Analysis.}  
% Let $n=|V|$, $m=|E|$, and let $K_s$ be the set of $s$-cliques. For a given $r < s$, define the $(r,s)$-degree of an $r$-clique $C_r$ as the number of $s$-cliques containing it, denoted $\deg_{r,s}(C_r)$. Then let
% \[
% \Sigma_{r,s} = \sum_{C_r} \deg_{r,s}(C_r) = \binom{s}{r} \cdot |K_s|, \qquad
% \Delta_{r,s} = \max_{C_r} \deg_{r,s}(C_r).
% \]

% Line 1 computes $s$-clique counts for all $r$-cliques using the CPI representation. Each $s$-clique contributes a constant amount of work along a root-clique path and is charged to all of its $\binom{s}{r}$ constituent $r$-cliques. Therefore, the time complexity is $O(\Sigma_{r,s})$ and the space cost is $O(n+\Sigma_{r,s})$. Line 2 builds a min heap $H_{\min}$ over all $r$-cliques, costing $O(N_r)$ time, where $N_r=\binom{n}{r}$ is the number of $r$-cliques.

% Each iteration extracts the $r$-clique of minimum count (Line 5), which costs $O(\log N_r)$. After extracting $S_{\min}$, the algorithm removes $S_{\min}$ from all root-clique paths in $\mathcal{T}$ that contain it and updates the counts of affected $r$-cliques (Lines 8-12). Across the whole execution, every $s$-clique is processed exactly once for each of its $\binom{s}{r}$ constituent $r$-cliques, so the total number of updates is $\Sigma_{r,s}$. Each update requires constant work on the CPI plus a heap decrease-key on $H_{\min}$, yielding $O(\log N_r)$ per update. Thus, the total cost of all updates is $O(\Sigma_{r,s}\log N_r)$.

% The \textsc{Batch Remove $r$-Cliques} procedure operates on a root-clique path $\mathcal{P}$ and a family of $r$-cliques to remove. Its recursive function \textsc{PathSplit} ensures that each affected $s$-clique on $\mathcal{P}$ is considered once, and each constituent $r$-clique is updated a bounded number of times. Therefore, its complexity is proportional to the number of affected $s$-cliques on $\mathcal{P}$, already accounted for in the $O(\Sigma_{r,s})$ update bound.

% Summing all contributions, the total running time of Algorithm~9 is
% \[
% O\big(\underbrace{\Sigma_{r,s}}_{\textsc{$s$-CliqueCount}} + \underbrace{N_r}_{\text{heap build}} +
% \underbrace{N_r \log N_r}_{\text{all extracts}} + \underbrace{\Sigma_{r,s}\log N_r}_{\text{all updates}}\big)
% = O\big((N_r + \Sigma_{r,s})\log N_r + \Sigma_{r,s}\big).
% \]
% For constant $r$ and $s$, this simplifies to $O\big((\binom{n}{r} + \binom{s}{r}|K_s|)\log \binom{n}{r}\big)$. The space complexity is dominated by the CPI representation and the heap, giving $O(n+\Sigma_{r,s})$ overall.

% \myaddRA{
\stitle{Complexity.} 
Let a clique path $\TPath = \{V_p(\TPath), V_h(\TPath)\}$.
The worst-case time complexity of Algorithm~\ref{alg:removeRClique} is $O(\prod_{C \in C_{\mathrm{rm}}} |C \cap V_p(\TPath)|)$, as it may 
recursively split for each clique in $C_{\mathrm{rm}}$. However, in practice, due to $s$-pruning and overlap among conflict cliques, the actual number of splits and the running time are significantly lower than this theoretical bound.
% }
